\section{Dataset}
\label{sec:dataset}

All data for this project is collected from free, publicly available
sources. Team-level NBA game logs and advanced statistics are pulled
from the \texttt{nba\_api} Python package, which exposes the
\href{https://stats.nba.com}{NBA.com} stats endpoints. Market-implied
probabilities are pulled from the Kalshi REST API at
\texttt{api.elections.kalshi.com/trade-api/v2}. The two pipelines are
implemented in \texttt{scripts/fetch\_data.py} and
\texttt{scripts/fetch\_kalshi\_data.py}; their outputs are merged on
game-level keys to produce the supervised training matrix used by every
model in this report.

\subsection{NBA Feature Engineering}

For each game $g_k$ we construct a $14$-dimensional pregame feature
vector $\tilde{\mathbf{x}}_k$ from team game logs. The features
characterize matchup state, recent form, and pace:
\begin{itemize}
    \setlength\itemsep{0pt}
    \item rest days since each team's previous game (home, away);
    \item season-to-date home/away split win percentage (home, away);
    \item rolling $5$-game and $10$-game win percentage (home, away);
    \item rolling $5$-game offensive rating, defensive rating, and pace
          of play (home, away).
\end{itemize}
Each feature is computed by aggregating each team's prior-games window
and is then \emph{temporally shifted} so that row $g_k$ contains only
information available strictly before $g_k$ tips off. This eliminates
same-game and future-game leakage and matches the conditions under
which the model would be queried in live deployment.

\subsection{Kalshi Market Data and Pre-Game Price Snapshots}

\textbf{Market scope.}
Kalshi's NBA series exposes contracts for many distinct outcomes,
including first-half winners, quarter winners, player-prop totals, and
series winners. Because our supervised target is full-game home-team
victory, we restrict our market matcher to event tickers prefixed with
\texttt{KXNBAGAME-}, which corresponds exclusively to the full-game
moneyline contract. All other contract families are filtered out at
ingestion in \texttt{src/data/kalshi\_fetcher.py}.

\textbf{Pre-tipoff price extraction.}
Kalshi's per-market bid/ask snapshot reflects the most recent quote,
which for settled markets converges on \$0 or \$1 and is therefore
useless as an estimate of pregame market belief. To recover a
pre-tipoff price for each matched market we query the
\texttt{/markets/candlesticks} endpoint at one-minute granularity. We
estimate tip-off as
\[
t_{\text{tip}} = t_{\text{close}} - 165\text{ min},
\]
where $t_{\text{close}}$ is the market's \texttt{close\_time} (trading
ceases at game end) and 165 minutes is the average NBA game length
including stoppages. From the candles in
$[t_{\text{tip}} - 3\text{h},\,t_{\text{tip}} - 5\text{min}]$ we select
the most recent candle with non-null \texttt{yes\_bid} and
\texttt{yes\_ask} closes, take the bid/ask midpoint as the YES price,
and flip to $1 - p_{\text{YES}}$ when the market's YES side
corresponds to the away team. The result is the
\texttt{kalshi\_price\_home\_win} column used as our market-implied
home-win probability.

\textbf{Coverage limitation.}
Of the $882$ KXNBAGAME markets matched to NBA games in our window,
$150$ ($17\%$) have qualifying pre-game candles. The remainder fall
back to the most recent available quote, which for thinly traded
markets can sit near \$0 or \$1; this artefact disproportionately
affects ROI estimates on the long tail of low-liquidity contracts and
is acknowledged as a limitation of our economic evaluation in
Section~\ref{sec:evaluation}.

\subsection{Merging and Final Dataset}

The NBA and Kalshi tables are joined on
\texttt{(GAME\_DATE, home\_TEAM\_ABBREVIATION,
away\_TEAM\_ABBREVIATION)}; rows lacking a matched Kalshi home-win
price are dropped. The resulting dataset spans 2025-04-15 through
2026-03-24 and contains $882$ games. We split the merged data
\emph{chronologically} 80/20: the earliest $705$ games form the
training set and the most recent $177$ games are held out as the test
set. No random shuffling is performed, so test-time evaluation always
sees games strictly later than any game seen in training, mirroring
real-time forecasting. Counts are summarized in Table~\ref{tab:dataset}.

\begin{table}[h]
\centering
\caption{Final moneyline dataset after Kalshi merge and chronological split.}
\label{tab:dataset}
\begin{tabular}{lr}
\toprule
\textbf{Field} & \textbf{Value} \\
\midrule
Total merged games           & 882 \\
Pregame features per game    & 14 \\
Train rows (earliest 80\%)   & 705 \\
Test rows (latest 20\%)      & 177 \\
Train date range             & 2025-04-15 -- 2026-02-21 \\
Test date range              & 2026-02-21 -- 2026-03-24 \\
KXNBAGAME markets matched    & 882 \\
Markets with pre-game candle & 150 (17\%) \\
\bottomrule
\end{tabular}
\end{table}

\section{Evaluation}
\label{sec:evaluation}

We evaluate each model along two complementary dimensions:
\textbf{predictive quality} on the held-out test set, and
\textbf{economic value} when the model's probabilities are used to
trade against Kalshi market prices. All experiments are reproducible
from the project repository; the training pipeline is implemented in
\texttt{scripts/train.py} and the simulated trading harness in
\texttt{scripts/backtest.py}, both of which read from
\texttt{config/config.yaml}.

\subsection{Experiment Settings}

All three classifiers are trained on the $14$-feature pregame matrix
defined in Section~\ref{sec:dataset} using the chronological 80/20
split (705 training games, 177 test games). Logistic regression
features are standardized with a \texttt{StandardScaler} fit on the
training fold; the tree-based models operate on raw feature scales.
Hyperparameters are summarized in Table~\ref{tab:hyperparams}; the
XGBoost values were selected via a small grid sweep on the training
fold (\texttt{scripts/tune\_xgboost.py}). Random seeds are fixed
at $42$ throughout to ensure repeatability.

\begin{table}[h]
\centering
\caption{Final hyperparameters for each model.}
\label{tab:hyperparams}
\begin{tabular}{ll}
\toprule
\textbf{Model} & \textbf{Hyperparameters} \\
\midrule
Logistic Regression & $C=1.0$, $\ell_2$ penalty, max\_iter $=1000$ \\
Random Forest       & $n_{\text{est}}=100$, depth $=15$, min\_split $=10$ \\
XGBoost             & $n_{\text{est}}=45$, depth $=6$, lr $=0.05$, \\
                    & subsample $=0.8$, colsample $=0.8$ \\
\bottomrule
\end{tabular}
\end{table}

\subsection{Predictive Quality}

We evaluate each model with four metrics on the test set
$\mathcal{G}_{\text{test}}$: classification accuracy, area under the
ROC curve (ROC AUC), log-loss
$\mathcal{L}_{\text{log}} = -\tfrac{1}{|\mathcal{G}_{\text{test}}|}
\sum_{g_k}[y_k \log \hat{p}_k + (1-y_k)\log(1-\hat{p}_k)]$, and the
Brier score
$\text{BS} = \tfrac{1}{|\mathcal{G}_{\text{test}}|}
\sum_{g_k}(\hat{p}_k - y_k)^2$. Log-loss and Brier score are
emphasized because identifying mispriced contracts requires
well-calibrated probability estimates, not merely correct
classifications. Results are reported in
Table~\ref{tab:predictive}.

\begin{table}[h]
\centering
\caption{Predictive quality on the held-out test set ($n=177$). Lower
is better for log-loss and Brier; higher is better for accuracy and
ROC AUC.}
\label{tab:predictive}
\begin{tabular}{lcccc}
\toprule
\textbf{Model} & \textbf{Acc.} & \textbf{ROC AUC} & \textbf{Log-loss} & \textbf{Brier} \\
\midrule
Logistic Regression & 0.7232 & 0.7603 & 0.5799 & 0.1968 \\
Random Forest       & 0.7062 & 0.7761 & 0.5741 & 0.1945 \\
XGBoost             & 0.6780 & 0.7368 & 0.6022 & 0.2067 \\
\bottomrule
\end{tabular}
\end{table}

Contrary to the prior expressed in our progress report and in much of
the tabular-prediction literature, the linear logistic baseline
attains the highest test accuracy ($72.3\%$), while the gradient-boosted
ensemble is in fact the weakest of the three on every metric we
report. Random Forest produces the best calibration (lowest Brier
score and highest ROC AUC). We attribute the lack of lift from
XGBoost to two factors: the modest size of the training fold ($705$
games) limits the data available to fit a high-capacity sequential
ensemble, and the $14$ engineered features encode signal that is
already largely linearly separable, leaving little nonlinear residual
for boosting to exploit.

\subsection{Economic Value: Simulated Trading}

\textbf{Strategy.}
For each test-set game we compute the model's edge over Kalshi on each
side, $\Delta_{\text{home}} = \hat{p}_k -
p^{\text{Kalshi}}_{\text{home},k}$ and analogously
$\Delta_{\text{away}}$. When the larger edge exceeds a threshold
$\tau$, a fixed-stake bet of \$10 is placed at the corresponding
Kalshi price. A correct bet pays
$10\cdot(1/p^{\text{Kalshi}}) - 10$; a wrong bet loses the stake.
Return on investment is reported as
$\text{ROI} = \text{net P\&L} / \text{total wagered}$.

\textbf{Clean evaluation subset.}
As noted in Section~\ref{sec:dataset}, only $150$ of the $177$ test
games carry a candle-derived pre-tipoff price; the remaining $27$
fall back to a most-recent quote that can be pinned near \$0 or \$1
in low-liquidity markets, which produces unrealistically large
simulated payouts on the few long-tail bets that resolve correctly.
To prevent this measurement artefact from dominating the headline
numbers, all ROI results in this subsection are computed on the
\emph{clean subset} of $n=150$ games for which a pre-game candle was
available. Numbers reported on the full $177$-game test set are
qualitatively similar but inflated in magnitude.

\textbf{Threshold sweep.}
Table~\ref{tab:sweep} reports XGBoost ROI across a range of edge
thresholds on the clean subset. The trend is unambiguous: tightening
$\tau$ does not yield a profitable subset of high-confidence bets.
ROI is between $-28\%$ and $-36\%$ across every threshold tested, and
the headline strategy at $\tau=0.05$ loses $\$428$ out of $\$1{,}190$
wagered. The strongest disagreements between model and market are
precisely the bets that lose the most money in absolute terms,
consistent with \emph{market efficiency} on Kalshi NBA moneyline
contracts within our test window: cases where our model deviates from
market consensus are not mispricings to be exploited, but cases where
the market incorporates information our feature set does not capture.
Figure~\ref{fig:roi-threshold} visualizes this together with bet count.

\begin{table}[h]
\centering
\caption{XGBoost ROI vs.\ edge threshold $\tau$ on the clean subset
($n=150$).}
\label{tab:sweep}
\begin{tabular}{cccr}
\toprule
$\tau$ & \textbf{Bets} & \textbf{Wagered (\$)} & \textbf{ROI (\%)} \\
\midrule
0.00 & 150 & 1{,}500 & $-28.6$ \\
0.02 & 137 & 1{,}370 & $-29.6$ \\
0.05 & 119 & 1{,}190 & $-36.0$ \\
0.08 & 105 & 1{,}050 & $-34.1$ \\
0.10 &  91 &   910   & $-32.7$ \\
0.15 &  65 &   650   & $-34.4$ \\
\bottomrule
\end{tabular}
\end{table}

\begin{figure}[h]
\centering
\includegraphics[width=0.85\columnwidth]{figures/roi_vs_threshold.png}
\caption{XGBoost ROI (left axis) and bet count (right axis) as a
function of the edge threshold $\tau$ on the clean subset. ROI is
negative and roughly flat in $\tau$, indicating no profitable
sub-population of high-confidence model picks.}
\label{fig:roi-threshold}
\end{figure}

\textbf{Strategy comparison.}
At our default threshold $\tau = 0.05$, we compare each learned model
against the two baselines specified in the proposal: a \emph{naive}
strategy that always bets the home team, and a \emph{kalshi-favorite}
strategy that bets whichever side Kalshi prices above $0.5$ (the
market's own pick). Results on the clean subset are reported in
Table~\ref{tab:strategy}.

\begin{table}[h]
\centering
\caption{Strategy comparison on the clean subset ($n=150$), \$10 flat
stake, $\tau=0.05$ for learned models. The kalshi-favorite baseline
is the only strategy with positive ROI.}
\label{tab:strategy}
\begin{tabular}{lrr}
\toprule
\textbf{Strategy} & \textbf{Bets} & \textbf{ROI (\%)} \\
\midrule
Kalshi-favorite (market pick) & 150 & $+6.3$  \\
Always-home (naive)           & 150 & $-2.7$  \\
Logistic Regression           & 121 & $-18.8$ \\
Random Forest                 & 120 & $-18.3$ \\
XGBoost                       & 119 & $-36.0$ \\
\bottomrule
\end{tabular}
\end{table}

Every learned model loses money over the test window, while
mechanically following Kalshi's own implied favorite returns
$+6.3\%$. Two interpretations are mutually compatible. First, Kalshi
prices are well-calibrated with a small bias toward favorites that a
disciplined bettor can capture without any modeling at all,
consistent with the prediction-market literature surveyed in
Section~2~\cite{Wolfers2004}. Second, our feature set is genuinely
weaker than the information embedded in the contract prices: rolling
team statistics ignore injuries, lineup changes, late-breaking news,
and matchup-specific style effects that traders on Kalshi routinely
incorporate. Reproducing the kalshi-favorite ROI with a \emph{learned}
model would require feature inputs at least as rich as the market's
own information set.

\subsection{Calibration}

Figure~\ref{fig:calibration} compares the reliability of the XGBoost
predicted probabilities against the Kalshi market probabilities on the
same test set. Kalshi tracks the diagonal more closely across the
probability range, while the model's curve sags below the diagonal in
the high-probability bins, confirming the Brier-score gap in
Table~\ref{tab:predictive} and explaining why the ROI gap is most
pronounced at large edge thresholds: the model is overconfident
exactly where it is wrong.

\begin{figure}[h]
\centering
\includegraphics[width=0.85\columnwidth]{figures/calibration.png}
\caption{Reliability diagram for the XGBoost model and the Kalshi
market on the held-out test set, with $10$ quantile-binned
probability buckets. Kalshi tracks the diagonal more closely,
indicating better calibration than the learned model.}
\label{fig:calibration}
\end{figure}

\subsection{Exploratory Observation: Reversed-Edge Strategy}
\label{sec:reverse}

The unambiguously negative ROI of the standard edge strategy
(Section~6.3) suggests an inverted experiment: instead of betting on
the side where the model's probability \emph{exceeds} the market, bet
on the side where it \emph{falls short}. Concretely, when
$\Delta_{\text{home}} > \tau$ we now place the stake on the away side
at price $1 - p^{\text{Kalshi}}_{\text{home}}$, and symmetrically for
$\Delta_{\text{away}}$. Identical bet count and stake schedule, with
the side-selection rule reversed.

Table~\ref{tab:reverse} reports the results for XGBoost (the model
with the largest model–market disagreement) across the same threshold
grid. Every threshold yields a positive simulated ROI; the strategy
returns $+12.0\%$ at $\tau=0$ and remains profitable through
$\tau=0.15$. At the headline threshold $\tau=0.05$, the inverted
strategy returns $+9.0\%$ on $\$1{,}190$ wagered—roughly 2.7 percentage
points above the kalshi-favorite baseline on the same set of games.

\begin{table}[h]
\centering
\caption{Reversed-edge strategy (XGBoost), \$10 flat stake, clean
subset. The strategy bets \emph{against} the model's directional
prediction whenever $|\Delta| > \tau$.}
\label{tab:reverse}
\begin{tabular}{cccr}
\toprule
$\tau$ & \textbf{Bets} & \textbf{Wagered (\$)} & \textbf{ROI (\%)} \\
\midrule
0.00 & 150 & 1{,}500 & $+12.0$ \\
0.02 & 137 & 1{,}370 & $+11.1$ \\
0.05 & 119 & 1{,}190 & $+9.0$  \\
0.08 & 105 & 1{,}050 & $+9.3$  \\
0.10 &  91 &   910   & $+4.4$  \\
0.15 &  65 &   650   & $+1.6$  \\
\bottomrule
\end{tabular}
\end{table}

We flag this as an \emph{exploratory} observation rather than a
deployable strategy for three reasons. First, the test window contains
only $n=150$ priced games over a 31-day stretch in 2026; the
$\sim$2.7-point lift over kalshi-favorite is plausibly within sampling
noise at that scale, and the strategy's apparent edge could vanish on
an out-of-sample test set. Second, treating systematic
anti-calibration as a profit signal is methodologically suspicious: it
asks the reader to trust that our model is consistently \emph{wrong}
in a stable direction, which is harder to justify a priori than a
correctly-signed signal of comparable strength. Third, the same
inversion run on Logistic and Random Forest yields only $+0.7\%$ and
$+1.4\%$ respectively at $\tau=0.05$, so the effect appears specific
to XGBoost rather than reflecting a property shared by all our
classifiers. We present the result as a hypothesis worth replicating
across a larger test horizon and a richer feature set before drawing
deployment-grade conclusions.

\textbf{Limitations.}
Two caveats temper the ROI numbers above. First, the headline
strategy comparison is computed on the clean subset of $150$ games;
the $27$ excluded games carry fallback quotes pinned near \$0 or \$1
that produce measurement noise rather than meaningful trading signal,
and re-introducing them inflates magnitudes without changing
directions (every learned model still loses money). Second, the test
window covers a single 31-day stretch in 2026; ROI estimates carry
meaningful sampling variance at $n=150$ games, and this evaluation
should be repeated across a longer horizon before treating the
negative-ROI conclusions as definitive.
